\section{Constitutional Governance and Co-Evolution (ARI-RQGM)}\label{sec:rqgm}

ARI's \term{institution} is the set of prompt-defined roles that propose
research directions, review artefacts, judge quality, and write the
manuscript. \term{ARI-RQGM} lets that institution evolve under a
constitution that does not. An epoch-based governance layer makes prompts
and components compete, attack, defend, and become adopted or retired,
while every institutional state change is validated by a deterministic,
non-model kernel against rule tables frozen in code. Governance is
selected before a run begins, and both the resolved mode and its
provenance are recorded in the checkpoint.
All changes commit at a boundary; T16 emergency quarantine forces an
immediate close/open boundary with a fresh institutional fingerprint.

\subsection{Motivation and lineage}\label{sec:rqgm:motivation}

The institutional roles are themselves prompt-defined artefacts, so
letting them evolve promises the same kind of improvement the search
loop extracts from experiments. Unconstrained self-modification,
however, is precisely what the G\"odel-machine literature warns
about: a component that can rewrite its own evaluator will eventually
reward itself. The Red Queen G\"odel Machine
\cite{iacob2026red} adds the co-evolutionary counterweight --- agents
improve \emph{because} adversaries and evaluators co-evolve against
them --- and ARI-RQGM takes its name and its central dynamic from that
line of work. But co-evolution alone does not bound the failure modes
that matter for a research system: score inflation, metric gaming,
evidence fabrication, and collusion between roles.

ARI-RQGM therefore places the co-evolutionary loop under a fixed
constitutional layer that never evolves. Its availability policy permits
candidate generation and local computation to continue after governance
failure only under a side-effect-free sandbox assumption. The kernel does
not enforce isolation of external APIs, devices, or production data;
irreversible operations need a separate fixed gate. State changes remain
fail-closed.

The determinism principle of \cref{sec:overview:principles} is narrower
than model reproducibility. Given the same saved advisory report,
candidate evaluations, registry, and configuration, transition resolution
and frontier repair are pure functions. The same committed event sequence
reconstructs the same projection. We do not claim that another model call
regenerates the same defence, adjudication proposal, or review; the three
boundaries are specified in \cref{sec:appendix:determinism}.

\subsection{Three layers, four facades}\label{sec:rqgm:layers}

Governance stratifies the system into three layers with strictly
decreasing authority to change things and strictly increasing freedom
to be changed.

\begin{description}
  \item[\textbf{Layer 0, constitutional (fixed).}] The deterministic
        kernel, the fixed result-verification path, and the
        append-only audit log. These never evolve; the rule tables
        they enforce live in code, pinned by a constitution hash
        (\cref{sec:rqgm:constitution}), and are registered in the
        component registry for provenance only.
  \item[\textbf{Layer 1, institutional.}] The evolvable service
        roles: reviewer, adversary, defender, artefact judge, router,
        utility policy, and the paper writer/reviewer when the archive
        phase is active. Their prompts or policy may be replaced through
        the lifecycle of \cref{sec:rqgm:prompts}, at epoch boundaries
        only. The auditor, evidence clerk, and governance judge are also
        registered and sanctionable, but v1 deliberately gives these
        judiciary actors no successor-generation path.
  \item[\textbf{Layer 2, meta.}] The agents that evolve Layer~1 ---
        the prompt and policy mutators, clean-room generator, replay
        selector, and failure-summary compressor --- themselves governed, with strictly
        less authority than the layer they modify: a meta agent can
        propose a replacement, never appoint one.
\end{description}

Four facades own the machinery. The \term{constitutional kernel}
validates records, hashes, capabilities, transitions, role
separation, erasure, audit-log integrity, clean-room inputs, and
context scope --- twelve closed entry points, none of which calls a
model or reads a clock. The \term{governance orchestrator} runs the
epoch-boundary audit (\cref{sec:rqgm:boundary}) and produces an
advisory report; it never writes a registry. The \term{registry
transition engine} is the sole writer of component and prompt
statuses, resolving changes against a fixed transition table. The
\term{frontier repair engine} restores the search frontier after a
retirement (\cref{sec:rqgm:boundary}). Around them sit the per-node
services: a proposal router, an adversarial round, a prompt-evolution
pipeline, a clean-room coordinator, and a budget manager that can
degrade any governance decision point without ever touching node
execution.

\subsection{Boot: mode resolution and founding registration}\label{sec:rqgm:boot}

\begin{figure}[t]
  \centering
  \resizebox{\columnwidth}{!}{% F22: full lifecycle of a governed (ari_rqgm) run, set in the ARI house
% design language (_style.tex). Three phase columns, F21-style: BOOT
% (lavender / ideation fill) — mode resolution, the one-time founding
% registration transaction (25 prompts + 12 components), and the first
% epoch freeze; GOVERNED SEARCH (peach / execution fill) — the per-node
% hooks (proposal round, ungoverned node execution, governance-level
% assignment, adversarial round) inside a dashed loop enclosure; EPOCH
% BOUNDARY (pale green / gate fill) — audit, kernel-validated transition,
% frontier repair + clean room, and the next freeze. The dotted margin
% path at the bottom is the reflux into the next epoch's search. No
% accent edge: this figure has no single emphasis role (F23 carries it).
%
% style.tikzstyles is loaded by the preamble / the standalone wrapper;
% the ari* house styles come from _style.tex. The build cwd differs per
% entry point, so resolve _style.tex from each known prefix.
\InputIfFileExists{../shared/figures/tikz/_style.tex}{}{%
  \InputIfFileExists{shared/figures/tikz/_style.tex}{}{%
    \InputIfFileExists{report/shared/figures/tikz/_style.tex}{}{%
      \InputIfFileExists{_style.tex}{}{%
        \errmessage{F22: cannot locate figures/tikz/_style.tex}}}}}
\begin{tikzpicture}[aricanvas]

  % --- column 1: boot (mode + founding registration) -------------------
  \node[ariboxwide, fill=phaseidea]           (mode)     {\figlabel{F22.mode}};
  \node[ariboxwide, fill=phaseidea, below=of mode] (founding) {\figlabel{F22.founding}};
  \node[ariboxwide, fill=phaseidea, below=of founding] (freezeo) {\figlabel{F22.freeze0}};
  \node[ariphase, above=0.4cm of mode]        (h1)       {\figlabel{F22.head.boot}};
  \pic at ([xshift=-2.5mm]h1.west) {ariglyph bulb};
  \draw[ariarrow] (mode)     -- (founding);
  \draw[ariarrow] (founding) -- (freezeo);

  % --- column 2: governed search (per-node hooks, loop-enclosed) -------
  \node[aribox exp, right=of mode]            (proposal) {\figlabel{F22.proposal}};
  \node[aribox exp, below=of proposal]        (execute)  {\figlabel{F22.execute}};
  \node[aribox exp, below=of execute]         (level)    {\figlabel{F22.level}};
  \node[aribox exp, below=of level]           (advers)   {\figlabel{F22.adversarial}};
  \node[ariloop, fit=(proposal)(execute)(level)(advers)] (loopgov) {};
  \node[ariphase, above=0.65cm of proposal]   (h2)       {\figlabel{F22.head.search}};
  \pic at ([xshift=-2.5mm]h2.west) {ariglyph flask};
  \node[arinote, anchor=south west] at ([xshift=6mm]loopgov.north) (ttl) {\figlabel{F22.loop}};
  \pic at ([xshift=-2.5mm]ttl.west) {ariglyph loop};
  \draw[ariarrow] (proposal) -- (execute);
  \draw[ariarrow] (execute)  -- (level);
  \draw[ariarrow] (level)    -- (advers);

  % --- column 3: epoch boundary ----------------------------------------
  \node[aribox gate, right=of proposal]       (audit)    {\figlabel{F22.audit}};
  \node[aribox gate, below=of audit]          (transit)  {\figlabel{F22.transition}};
  \node[aribox gate, below=of transit]        (repair)   {\figlabel{F22.repair}};
  \node[aribox idea, below=of repair]         (freezen)  {\figlabel{F22.freeze1}};
  \node[ariphase, above=0.4cm of audit]       (h3)       {\figlabel{F22.head.boundary}};
  \pic at ([xshift=-2.5mm]h3.west) {ariglyph gate};
  \draw[ariarrow] (audit)   -- (transit);
  \draw[ariarrow] (transit) -- (repair);
  \draw[ariarrow] (repair)  -- (freezen);

  % --- handoffs between the phase columns -------------------------------
  \draw[ariarrow, rounded corners=5pt]
    (freezeo.east) -- ([xshift=8mm]freezeo.east) |- (proposal.west);
  \draw[ariarrow, rounded corners=5pt]
    (advers.east) -- ([xshift=8mm]advers.east)
    |- node[arinote, pos=0.3] {\figlabel{F22.note.trigger}} (audit.west);

  % --- reflux: the new epoch keeps searching ----------------------------
  \draw[arifeedback]
    (freezen.south) -- ([yshift=-7mm]freezen.south) -| (loopgov.south);
  \node[arinote] at ([yshift=-7mm]loopgov.south east) {\figlabel{F22.note.next}};

\end{tikzpicture}
}
  \caption{Lifecycle of a governed run. Boot resolves the execution
    mode, executes the one-time founding registration transaction (32
    prompt-or-policy records, 20 components), and freezes the founding active set as
    epoch~000. During the epoch, per-node hooks record proposals,
    assign a governance level (L0 fixed to L3 adjudicated), and run
    an event-driven adversarial round; node execution itself is never
    vetoed. When the epoch's node-count trigger fires, the boundary
    runs the advisory audit, the kernel-validated transition, and
    frontier repair with clean-room regeneration, then freezes the
    next epoch's active set and the search continues.}
  \label{fig:rqgm-lifecycle}
\end{figure}

A governed run boots in three steps (\cref{fig:rqgm-lifecycle}).
First, \term{mode resolution}: the mode is a two-key interlock --- a
mode selector and a redundant enable flag must agree, with
environment overrides validated before assignment --- and any
mismatch warns and leaves the constitutional layer inactive rather than
guessing. The resolved mode and its provenance (configuration,
environment, or resume) are recorded in the checkpoint.

Second, the runtime attaches governance as best-effort hooks around the
research executor. The executor retains authority over scientific work;
the hooks observe typed proposals, evaluations, and artefacts without
acquiring authority to veto their execution.

Third --- on a fresh governed checkpoint only --- the \term{founding
registration} transaction populates the institution. From two tables
frozen in code, one prepare event is followed by 32 prompt
registrations and 20 component registrations, then a single commit.
The 20 founding components are the research generator, seven exploration-adversary variants,
artefact judge, auditor, clean-room generator, defender, evidence clerk,
failure-summary compressor, governance judge, policy mutator, prompt
mutator, proposal router, replay selector, and utility policy. The 32
founding prompts comprise 31 committed templates plus the configuration-
derived utility-policy prompt. The paper-archive phase registers three
further prompts and components --- paper writer, paper reviewer, and the
paper-self-preference adversary --- for totals of 35 and 23. Each prompt
is registered with the hash of its template bytes. The transaction then
freezes the founding active
set --- active components, active prompt hashes, and the utility
policy --- into the immutable state of epoch~000. Its policy fingerprint
also binds the resolved governance settings, while its separately reported
execution fingerprint binds declared model, decoding, tool, environment, and
data identities; missing external revisions remain marked unresolved. The
composite deliberately excludes timestamps. Founding is write-once: on resume the
registries are reconstructed by replaying the event log, a committed
registration is never re-run, and a founding transaction interrupted
before its commit is invisible and simply re-executed.

\subsection{Per-node governance}\label{sec:rqgm:node}

During an epoch the research executor continues to operate while
governance observes it through three best-effort hooks.

\paragraph{Proposals.}
When the expansion step generates candidate directions, the proposal
router archives a full \term{proposal record} --- the complete
generation context, on the principle of storing everything --- while
the search itself receives only a capped \term{proposal summary
view}. The summary is the \emph{only} surface the search sees, so
proposal generators can grow richer without widening the interface to
the search loop, and the archived records remain available to audits
and to replay.

\paragraph{Governance levels.}
Each evaluated node is assigned a governance level by a pure trigger
ladder --- from L0 (fixed checks only) to L3 (fully adjudicated) ---
based on deterministic signals: membership in the frontier's top
scorers, an anomalous score jump relative to the parent, a
paper-candidacy flag, a novelty claim, or low evaluator confidence.
The level decides how much governance attention the node receives,
never whether it runs. Every assignment is appended to the audit log,
and the budget manager may cap the effective level when a per-epoch
governance budget runs low --- degrading governance, not the node.
Continuation after governance failure is safe only for side-effect-free
candidate generation and local computation; this is an operating
condition, not a kernel invariant.

\paragraph{Adversarial rounds.}
Nodes that trip the same trigger family also receive an
\term{adversarial round}: one of seven exploration adversary types ---
overclaiming, evidence gaps, metric gaming, reproducibility,
prior-art, cost explosion, and prompt injection --- attacks the
node's \emph{artefacts}; the paper phase adds an eighth,
paper-self-preference adversary. Adversaries never attack components;
the raw-attack record has no component target by construction. A raw attack is
score-inert: only after the defender responds and the artifact judge
upholds the attack does a \term{validated attack record} exist, and
only validated attacks feed a bounded utility penalty under the
epoch-frozen penalty policy, with the pre-penalty score preserved in
additive metric keys. The judge-authored validated record may then bind
an implicated role to an epoch-frozen incumbent component for
accountability (\cref{sec:rqgm:targetbinding}). A judge failure defaults
to \emph{invalid} --- no penalty. This avoids an unverified sanction but
is vulnerable to denial of the adjudication path; fixed adjudication and
independent retry are not implemented. Validated attacks also enter the
\term{adversarial replay
pool}, the corpus that later tests candidate prompts
(\cref{sec:rqgm:prompts}) --- the Red-Queen coupling: today's
validated attacks are tomorrow's entrance exam. The round is
idempotent per node --- a marker in the case log ensures a resumed
run never re-attacks a node it already processed.

\subsection{The epoch boundary}\label{sec:rqgm:boundary}

\begin{figure}[t]
  \centering
  \resizebox{\columnwidth}{!}{% F23: the epoch-boundary governance pipeline, set in the ARI house design
% language (_style.tex). Left column: the nine-step advisory audit
% (observe .. self-audit, numbered badges; the dashed loop enclosure marks
% it read-only w.r.t. registries). Right column: the report, the
% transition engine's resolution, the constitutional kernel's validation
% of every edge against the fixed T1-T21 table, the four-event boundary
% transaction, and frontier repair. The single accent edge (the rationed
% emphasis role of this figure) is the kernel's fail-closed block of an
% illegal transition.
%
% style.tikzstyles is loaded by the preamble / the standalone wrapper;
% the ari* house styles come from _style.tex. The build cwd differs per
% entry point, so resolve _style.tex from each known prefix.
\InputIfFileExists{../shared/figures/tikz/_style.tex}{}{%
  \InputIfFileExists{shared/figures/tikz/_style.tex}{}{%
    \InputIfFileExists{report/shared/figures/tikz/_style.tex}{}{%
      \InputIfFileExists{_style.tex}{}{%
        \errmessage{F23: cannot locate figures/tikz/_style.tex}}}}}
\begin{tikzpicture}[aricanvas]

  % --- column 1: the nine-step epoch audit (advisory) -------------------
  \node[aribox exp]                        (observe)  {\figlabel{F23.observe}};
  \node[aribox exp, below=of observe]      (reliab)   {\figlabel{F23.reliability}};
  \node[aribox exp, below=of reliab]       (evidence) {\figlabel{F23.evidence}};
  \node[aribox exp, below=of evidence]     (prosec)   {\figlabel{F23.prosecute}};
  \node[aribox exp, below=of prosec]       (defense)  {\figlabel{F23.defense}};
  \node[aribox exp, below=of defense]      (adjud)    {\figlabel{F23.adjudicate}};
  \node[aribox exp, below=of adjud]        (pool)     {\figlabel{F23.replaypool}};
  \node[aribox exp, below=of pool]         (selfaud)  {\figlabel{F23.selfaudit}};
  \node[ariloop, fit=(observe)(reliab)(evidence)(prosec)(defense)(adjud)(pool)(selfaud)] (loopaud) {};
  \node[ariphase, above=0.65cm of observe] (h1)       {\figlabel{F23.head.audit}};
  \pic at ([xshift=-2.5mm]h1.west) {ariglyph flask};
  \node[arinote, anchor=north] at ([yshift=-2.5mm]loopaud.south) (ttl) {\figlabel{F23.note.advisory}};
  \draw[ariarrow] (observe)  -- (reliab);
  \draw[ariarrow] (reliab)   -- (evidence);
  \draw[ariarrow] (evidence) -- (prosec);
  \draw[ariarrow] (prosec)   -- (defense);
  \draw[ariarrow] (defense)  -- (adjud);
  \draw[ariarrow] (adjud)    -- (pool);
  \draw[ariarrow] (pool)     -- (selfaud);

  % Step badges 1..8 (the report itself is step 9).
  \node[aribadge] at (observe.north west)  {1};
  \node[aribadge] at (reliab.north west)   {2};
  \node[aribadge] at (evidence.north west) {3};
  \node[aribadge] at (prosec.north west)   {4};
  \node[aribadge] at (defense.north west)  {5};
  \node[aribadge] at (adjud.north west)    {6};
  \node[aribadge] at (pool.north west)     {7};
  \node[aribadge] at (selfaud.north west)  {8};

  % --- column 2: report -> transition -> kernel -> commit -> repair -----
  \node[aribox paper, right=2.4cm of observe] (report) {\figlabel{F23.report}};
  \node[aribox idea,  below=of report]        (resolve){\figlabel{F23.resolve}};
  \node[aribox gate,  below=of resolve]       (kernel) {\figlabel{F23.kernel}};
  \node[ariboxwide, fill=phasepaper, below=of kernel] (txn) {\figlabel{F23.txn}};
  \node[aribox exp,   below=of txn]           (repair) {\figlabel{F23.repair}};
  \node[arinote, below=0.2cm of repair]       (retn)   {\figlabel{F23.note.retire}};
  \node[ariphase, above=0.4cm of report]      (h2)     {\figlabel{F23.head.transition}};
  \pic at ([xshift=-2.5mm]h2.west) {ariglyph gate};
  \node[aribadge] at (report.north west) {9};
  \draw[ariarrow] (report)  -- (resolve);
  \draw[ariarrow] (resolve) -- (kernel);
  \draw[ariarrow] (kernel)  -- (txn);
  \draw[ariarrow] (txn)     -- (repair);

  % --- the audit hands its report across the column gap -----------------
  \draw[ariarrow, rounded corners=5pt]
    (selfaud.east) -- ([xshift=10mm]selfaud.east) |- (report.west);

  % --- the one accent role: an illegal edge is blocked fail-closed ------
  \node[ariboxwide, right=2.2cm of kernel] (blocked) {\figlabel{F23.blocked}};
  \draw[ariarrow accent] (kernel) -- (blocked);

\end{tikzpicture}
}
  \caption{The epoch-boundary pipeline. The nine-step audit is
    advisory and read-only with respect to registries; its report
    feeds the registry transition engine, whose resolved transition
    is validated edge-by-edge by the constitutional kernel against
    the fixed T1--T21 table before the four-event boundary
    transaction commits. An illegal edge is blocked fail-closed ---
    the incumbent active set keeps serving and the run continues.
    Frontier repair runs only when the committed transition retired a
    component or prompt.}
  \label{fig:rqgm-boundary}
\end{figure}

An epoch freezes the institution; the boundary is the only window in
which the institution may change. The boundary trigger is node count
(ten new nodes per epoch by default), checked on the main thread at
the head of each outer search iteration --- never with a node in
flight --- and once more at end of run, so an epoch whose final node
arrives in the last iteration still closes transactionally; a
trailing epoch below the trigger simply stays open. The end-of-run
check exists because live runs exposed the gap: a boundary reached
exactly at the last node would otherwise never fire, leaving the
epoch's audit unwritten.

\paragraph{The audit.}
The closing epoch is audited first (\cref{fig:rqgm-boundary}), by a
nine-step pipeline with a total deterministic fallback at every
LLM-consuming step, so an audit without a model still completes.
Observations are collected from the epoch's records; a reliability
assessment scores each governed component; the \term{evidence clerk}
--- and only the clerk --- assembles evidence bundles, excluding
same-role accusations, which are inadmissible by constitution;
prosecution follows, where only the \term{auditor} may file an
impeachment motion against a per-epoch motion quota; internal
\texttt{posted}/\texttt{refunded}/\texttt{forfeited} counters account for
that finite allowance but transfer no value and are a rate limit rather
than an economic bond; a defense is generated for every motion; an
adjudication panel rules; the replay pool is updated; and a
\term{governance self-audit} re-validates the governance actors' own
records through the kernel. These judiciary actors are founding
components and can themselves be sanctioned; when the governance judge
is the target of a motion it recuses and the deterministic board fallback
decides the case. The result is a \term{governance report}
--- reliability, motions, defenses, adjudications, candidate
evaluations, recommendations, motion-quota accounting --- appended to the
hash-chained audit log. The report is advisory input to the next
stage; the orchestrator that produced it cannot write a registry.

\paragraph{The transition.}
The registry transition engine resolves the report's recommendations
and the candidate evaluations into one epoch transition --- adoptions,
sanctions, retirements --- as a pure function of the frozen epoch
state, the report, the registries, and the configured thresholds:
byte-identical on re-run. Components and prompts move through ten
statuses (candidate, validated, shadow, probationary active, active,
warning, probation, quarantine, retired, banned) along a fixed table
of twenty-one rules, T1--T21. The kernel validates every edge of the
resolved transition: an edge outside the table is the named violation
\texttt{CK-REG-001}; a boundary-only edge stamped mid-epoch without
the emergency shape is \texttt{CK-REG-002}; a transition not authored
by the transition engine --- a judge writing the registry, say --- is
\texttt{CK-ROL-901}. A blocked transition aborts fail-closed with the
incumbent active set still serving; governance is suspended, the run
never is. A valid transition commits through a four-event boundary
transaction --- prepare, close the old epoch, open the new one,
commit --- so a crash between any two events leaves a prefix that
resume discards and re-runs deterministically, and a duplicate commit
of the same transition is a guarded no-op. \term{Emergency quarantine}
(rule T16) handles a kernel-critical violation without weakening epoch
identity. It force-closes the current epoch and commits the quarantine,
close, and opening of a freshly fingerprinted epoch in one transaction.
The prior utility policy is preserved and no replacement is adopted in
that emergency transaction. Nodes on the two sides therefore carry
different epoch identities; \cref{sec:appendix:transitions} gives the
complete component-lifecycle table.

\paragraph{Frontier repair.}
If the committed transition retired a component or prompt, the run's
past contains records produced under an institution now judged
defective, and the frontier must not keep paying interest on them.
The repair engine starts from records stamped with the retired prompt hash
and follows reverse provenance references breadth-first. Only
proposal-to-proposal inspiration is declared non-material; unknown edges
are material. The default depth cap is eight, with any still-reachable
records conservatively invalidated. It marks
the affected nodes stale via additive flags (stale, valid-for-frontier,
and the erasure reason) persisted through the tree file. Erasure is
\emph{logical only}: nothing is physically deleted, so the history
remains auditable while every consumer that promotes a node ---
frontier scoring and expansion, best-node selection, and the paper's
evidence grounding --- reads the flags and ignores it; when every
candidate is erased, selection reports no winner rather than falling
back to contaminated evidence. Repair is
role-aware and has two explicit recomputation paths. Utilities that
merely \emph{consumed} a retired reviewer's scored evidence are
recomputed from surviving inputs under the \emph{original} epoch's
frozen weights; records the retired role itself produced are
invalidated outright. By contrast, when the \emph{utility policy}
retires, every stamped node is re-scored under the successor policy
by recomposing its stored per-axis raw scores, and the repair event
lists those nodes in \texttt{policy\_rescored\_node\_ids}. Missing
raw axes fail closed to invalidation. This is an operational definition,
not an evaluated causal guarantee; semantic precision and over-removal
are unknown. If repair itself fails validation, the engine
degrades conservatively, ultimately to a drain-only mode in which the
run finishes pending work but expands no further --- never a crash.

\subsection{The evolving score}\label{sec:rqgm:utility}

Every subsection so far has let the institution evolve while treating one
thing as fixed: the \emph{score} itself --- the utility policy that ranks
the frontier, built from an axis basis and its weights, a composite rule,
a frontier-scoring rule, a depth penalty, and an exploration constant.
Freezing the score for the life of a run is the conservative choice, but
it is not the one the Red Queen line of work \cite{iacob2026red} actually
makes: its defining claim is that the ordering objective itself is
rewritten at each epoch boundary. Governed utility evolution closes that
gap. The utility policy becomes an evolvable role in its own right, with
one difference from the prompt roles of \cref{sec:rqgm:prompts} --- its
incumbent is a policy document, not a prompt template. It is frozen
within an epoch, so the frontier is ranked by one fixed rule while an
epoch runs, and it is rewritable only at the boundary, by the same
transition engine that adopts prompts.

A dedicated proposer mints candidate policies and nothing else. Its four
default mutations are pure arithmetic over the boundary's
already-abstract evidence --- no model, no clock, no randomness --- so
the default rewrite is fully reproducible, and it is forbidden from
reading the frontier's current scores: a score tuned to flatter the nodes
it has already produced is exactly the self-referential collusion the
co-evolutionary design exists to prevent. The proposer holds no registry
authority, and it is itself a registered, sanctionable component --- there
is no absolute ruler over the score any more than over the roles it
grades.

Adoption rides one new edge in the transition table. A successor policy
that has cleared validation and shadow service and scores at least as
well as the incumbent on the frozen replay board \emph{supersedes} it:
rule~T20, the sole active-to-retired edge, and the one retirement that
does not stage through quarantine. The kernel permits it for the
utility-policy role alone --- every behavioural component keeps the
conservative, sanction-only replacement model, and for them a direct
retirement is still rejected. Supersession retires the healthy incumbent
under its \emph{old} policy hash, and because every node now carries the
hash of the policy that scored it, frontier repair
(\cref{sec:rqgm:boundary}) selects every node ranked under the retired
policy, not merely the subset that happened to be penalised, and
re-composes its stored raw axis scores under the successor. Only a
node without recoverable axes is invalidated fail-closed.
Adding the edge re-pinned the constitution hash, as an amendment must
(\cref{sec:rqgm:constitution}).

One limit is worth stating as plainly as the mechanism. When the score's
axes are derived dynamically and no fixed per-axis basis is pinned, the
``scores at least as well'' gate has no stable ordering to compare
against, so supersession reduces to legality and non-degeneracy rather
than a demonstrated improvement. A gate that can certify a successor as
\emph{strictly} better needs an anchored basis --- ground truth to score
against --- which the paper pipeline of \cref{sec:rqgm:paper} supplies
but the exploration phase, lacking any external oracle over research
quality, does not. This is the honest Red-Queen posture: with no ground
truth, the criterion rolls within the space of legal policies rather than
climbing a fixed hill, and we do not present it as more than that.

\subsection{Binding an attack to an accountable component}\label{sec:rqgm:targetbinding}

The adversarial round of \cref{sec:rqgm:node} produces, after
adjudication, a validated attack record. It first names implicated
\emph{roles}; sanction additionally requires resolution to the component
that produced the attacked artefact under that epoch.

Each judge-authored record carries, whenever one is resolvable, a single
registry-resolvable target: the incumbent that made the decision
the attack invalidates, resolved when the record is minted against the
epoch-\emph{frozen} active set --- never a live registry read, which
after a boundary adoption would impeach a successor for its predecessor's
defect. Role separation is preserved exactly: the adversary attacks
artefacts and only artefacts, and it is the judge's verdict, not the
accusation, that makes an observation accountable, so no actor both
accuses and binds; a record that would bind its own adjudicator drops the
binding and says nothing rather than raise. When no target resolves, the
field is absent and no sanction follows. If multiple roles are
implicated, the system emits one record per role referencing the same
adjudication; the one-target constraint is an audit format, not a causal
guarantee.

The research generator is now a founding component. Before a generated
node is recorded or attacked, the runtime stamps its producer component,
prompt hash, and epoch exactly once from the frozen active set. Each of the
seven exploration attacks binds that generator only when all available
provenance agrees; legacy, missing, or mismatched provenance remains
targetless, preventing a successor from inheriting a predecessor's defect.
For the paper self-preference attack of \cref{sec:rqgm:paper}, one round
can bind the manuscript reviewer, the paper writer, or both registered
components. The fixed mapping's attribution accuracy and its treatment of
tool failures, data contamination, shared prompts, and genuinely joint
causes have not been evaluated.

\subsection{Prompt evolution, clean room, and the meta tier}\label{sec:rqgm:prompts}

\begin{figure}[t]
  \centering
  \resizebox{\columnwidth}{!}{% F24: the governed prompt lifecycle as a state machine, set in the ARI
% house design language (_style.tex). Column 1 (lavender / ideation fill)
% is the evolution side: the meta-tier clean-room generator and prompt
% mutator emit candidates, which climb the six-stage validation ladder
% (candidate -> validated -> shadow). Column 2 (peach / execution fill)
% is service: probationary_active -> active with the warning/probation
% sanction ladder below and the rehabilitation back-edge. Column 3
% (white) is the terminal sanction chain quarantine -> retired -> banned.
% The single accent edge (the rationed emphasis role of this figure) is
% T16 emergency quarantine, which forces an immediate close/open boundary.
% Dotted margin paths: exoneration re-entry (T18) and the
% clean-room regeneration loop from a retirement back to a fresh
% candidate. Rejection edges (candidate/shadow -> retired, T2/T5) are
% omitted from the geometry and stated in the caption.
%
% style.tikzstyles is loaded by the preamble / the standalone wrapper;
% the ari* house styles come from _style.tex. The build cwd differs per
% entry point, so resolve _style.tex from each known prefix.
\InputIfFileExists{../shared/figures/tikz/_style.tex}{}{%
  \InputIfFileExists{shared/figures/tikz/_style.tex}{}{%
    \InputIfFileExists{report/shared/figures/tikz/_style.tex}{}{%
      \InputIfFileExists{_style.tex}{}{%
        \errmessage{F24: cannot locate figures/tikz/_style.tex}}}}}
\begin{tikzpicture}[aricanvas]

  % --- column 1: evolution (meta tier emits, candidates climb) ----------
  \node[aribox idea]                       (cleanroom) {\figlabel{F24.cleanroom}};
  \node[aribox idea, below=of cleanroom]   (cand)      {\figlabel{F24.candidate}};
  \node[aribox idea, below=of cand]        (valid)     {\figlabel{F24.validated}};
  \node[aribox idea, below=of valid]       (shadow)    {\figlabel{F24.shadow}};
  \draw[ariarrow] (cleanroom) -- node[arinote] {\figlabel{F24.note.emit}} (cand);
  \draw[ariarrow] (cand)  -- (valid);
  \draw[ariarrow] (valid) -- (shadow);

  % --- column 2: service (adoption at the boundary; sanction ladder) ----
  \node[aribox exp, right=of cleanroom]    (probact)   {\figlabel{F24.probact}};
  \node[aribox exp, below=of probact]      (active)    {\figlabel{F24.active}};
  \node[aribox exp, below=of active]       (warning)   {\figlabel{F24.warning}};
  \node[aribox exp, below=of warning]      (probation) {\figlabel{F24.probation}};
  \draw[ariarrow] (probact) -- (active);
  \draw[ariarrow] (active)  -- (warning);
  \draw[ariarrow] (warning) -- (probation);
  \draw[ariarrow muted] (probation.west)
    to[bend left=28] node[arinote, pos=0.3] {\figlabel{F24.note.rehab}}
    (active.west);

  % adoption handoff: shadow (col 1) -> probationary_active (col 2)
  \draw[ariarrow, rounded corners=5pt]
    (shadow.east) -- ([xshift=8mm]shadow.east)
    |- node[arinote, pos=0.28] {\figlabel{F24.note.adopt}} (probact.west);

  % --- column 3: terminal sanctions --------------------------------------
  \node[aribox, right=of warning]          (quarantine) {\figlabel{F24.quarantine}};
  \node[aribox, below=of quarantine]       (retired)    {\figlabel{F24.retired}};
  \node[aribox, below=of retired]          (banned)     {\figlabel{F24.banned}};
  \draw[ariarrow] (quarantine) -- (retired);
  \draw[ariarrow] (retired)    -- (banned);

  % the one accent role: the forced emergency boundary edge (T16)
  \draw[ariarrow accent] (warning)
    -- node[arinote] {\figlabel{F24.note.emergency}} (quarantine);

  % exoneration re-entry (T18): back under probation
  \draw[arifeedback] (quarantine.north)
    |- node[arinote, pos=0.75] {\figlabel{F24.note.exon}} (probact.east);

  % clean-room regeneration: a retirement seeds a fresh candidate
  \draw[arifeedback]
    (retired.east) -- ([xshift=8mm]retired.east)
    |- ([yshift=7mm]cleanroom.north) -- (cleanroom.north);
  \node[arinote, above=0.25cm of probact] (reentry) {\figlabel{F24.note.reenter}};

\end{tikzpicture}
}
  \caption{The governed prompt lifecycle. Meta-tier agents (the
    prompt mutator; after a retirement, the clean-room generator)
    emit candidates only. A candidate climbs the six-stage validation
    ladder (static, constitutional, schema dry-run, replay against
    the adversarial pool, anchor, shadow) and can be adopted only by
    the transition engine at an epoch boundary, entering service
    under probation. Sanctions descend through warning and probation;
    the accented edge is T16 emergency quarantine, which forces an
    immediate close/open boundary. A retired prompt's text is sealed; its role is
    served by the baseline fallback until a clean-room replacement
    re-enters the lifecycle as a fresh candidate. Rejection edges
    from candidate and shadow to retired (T2, T5) are omitted for
    clarity.}
  \label{fig:prompt-lifecycle}
\end{figure}

Institutional prompts evolve through the lifecycle of
\cref{fig:prompt-lifecycle}. The \term{prompt mutator} emits
candidate templates and nothing else --- it holds no registry
authority, and there is deliberately no edge from candidate to
active, so no output of a model can go straight into service. A
candidate must clear six stages: \emph{static validation}
(deterministic well-formedness and no in-place mutation of the
incumbent), \emph{constitutional validation} (kernel checks on the
candidate record), a \emph{schema dry-run} (one injected call must
produce a reply parseable under the role's declared output schema),
\emph{replay evaluation} against the adversarial replay pool ---
candidates are graded on the very attacks that defeated their
predecessors --- an \emph{anchor evaluation} on a fixed task set
where ties favour the incumbent, and finally \emph{shadow} service,
in which the candidate runs side by side with the incumbent on
sampled traffic with observation-only effect. Only then can the
transition engine adopt it at a boundary, and adoption enters service
as probationary, not active.

Retirement triggers the strictest path in the system. A retired
prompt's text is sealed: registry views stub it out, and a dedicated
access guard denies it at the capability layer, everywhere, forever.
The replacement is produced in a \term{clean room}: the fixed tier
assembles a closed input bundle --- committed catalogues and abstract
failure summaries, never registry prompt text --- which the kernel
screens before generation (violations are the \texttt{CK-CLN}
family, including a word-shingle similarity screen against the
retired text at admission), and generation is one-shot, with no
tool-bearing loop that could go looking for the sealed text. The
regenerated candidate re-enters the lifecycle at the bottom rung;
meanwhile the retired role's slot is served by the built-in baseline,
so no governed role is ever vacant. The screen is lexical: semantic
paraphrase leakage and false rejection of legitimate role vocabulary have
not been measured. Clean-room replacement is therefore a contamination
reduction mechanism, not a guarantee of semantic non-inheritance.

The meta tier is governed more tightly than the layer it modifies.
Meta agents run in a read-only sandbox behind an allowlisted tool
proxy whose only write surface is a synthesized submission tool;
their capability flags are deny-by-default with the dangerous flags
constant-false; the kernel checks \term{authority non-expansion} as
subset arithmetic --- a candidate cannot declare more authority than
the serving incumbent it would replace (or a fixed role matrix when no
incumbent exists); and every meta output enters the lifecycle
at candidate status. The meta tier can propose; it cannot activate.
Nor is it failure-gated or silent: the mutator mints candidates at
every epoch boundary --- calm epochs included --- under per-role and
per-epoch budget caps, and each boundary's meta step records an
audit-visible summary of its outcome even when it proposes nothing.

\subsection{The constitution and its amendment}\label{sec:rqgm:constitution}

\begin{table}[t]
  \centering
  \small
  \begin{tabularx}{\columnwidth}{@{}lX@{}}
    \toprule
    \textbf{Family} & \textbf{What it guards} \\
    \midrule
    \texttt{CK-SCH} & Record schema: malformed governance records
      block; malformed node-side records warn. \\
    \texttt{CK-HSH} & Hash integrity: records, registered prompt
      hashes, artefact digests. \\
    \texttt{CK-ACC} & Capability: the deny-by-default authority
      matrix per role and tier. \\
    \texttt{CK-ROL} & Role separation; \texttt{CK-ROL-901} names the
      canonical violation --- anyone but the transition engine
      writing the registry. \\
    \texttt{CK-REG} & Transition legality against T1--T21:
      \texttt{001} an edge outside the table, \texttt{002} a
      boundary-only edge attempted mid-epoch. \\
    \texttt{CK-EPO} & Epoch invariance: the active set may not drift
      inside one fingerprinted epoch. \\
    \texttt{CK-ERA} & Selective erasure: stale records must not
      reach frontier scoring or best-node selection. \\
    \texttt{CK-AUD} & Audit-log integrity: the event hash chain and
      its monotonicity. \\
    \texttt{CK-CLN} & Clean-room hygiene: bundle screening and
      contamination against sealed text. \\
    \texttt{CK-CTX} & Context scope: role-scoped views may not leak
      beyond their role. \\
    \bottomrule
  \end{tabularx}
  \caption{Violation-code families of the constitutional kernel.
    Codes are stable identifiers, never renumbered; each validator
    reports the specific code it fired, and the enforcement adapter
    decides warn versus block by context --- state changes block,
    node-side findings warn.}
  \label{tab:ck-codes}
\end{table}

The constitution is code, not configuration. The transition table,
the capability matrix, the severity and erasure rules, and the
clean-room screens live in source, and a single \term{constitution
hash} derived from those tables is pinned by a guard test. Editing
any rule changes the hash and fails the guard, so a constitutional
change cannot happen as a side effect: the author must re-pin the
hash explicitly, alongside a dated rationale --- a deliberate,
reviewable act, closer to an amendment procedure than to a
configuration change. One amendment of this kind has occurred to
date: a meta-tier role present in the role vocabulary was missing
from the capability matrix, so every capability lookup for it was
refused; the amendment added the role with the same minimal grants as
its narrowest peer and re-pinned the hash with the rationale in
place. Numeric thresholds (budgets, trigger counts) are configurable;
the table topology is not. A human-readable rendering of the
constitution is copied into every governed checkpoint and its hash
recorded in the run's metadata --- editing the copy changes nothing,
by design. \Cref{tab:ck-codes} lists the violation-code families the
kernel reports.

\subsection{Governing the paper pipeline}\label{sec:rqgm:paper}

The same constitutional machinery also governs manuscript writing under
an independent paper-phase mode. The \emph{archive} mode places writing
under epoch-bound governance; leaving the mode inactive constructs no
paper-governance object. Like the run-level choice of
\cref{sec:rqgm:boot}, it is a two-key interlock whose selector and
redundant enable flag must agree.

The archive mode realises, for paper writing, the four commitments that
define the method. \emph{First, a draft archive under a rewritable score.}
The mode creates several seed drafts in distinct framings, iteratively
refines selected candidates, and ranks the eligible archive by the
governed reviewer's composite plus a diversity bonus for under-sampled
framings. A total-candidate budget bounds the work. Reviewer utility
co-evolves across epochs under governed utility evolution
(\cref{sec:rqgm:utility}), with repair re-scoring drafts from stored raw
axes under a successor policy.
\emph{Second,
adversaries attack evaluations, not only artefacts.} An eighth adversary
type attacks the reviewer's over-acceptance: an AI-authored draft the
incumbent reviewer scored highly that the human-labelled anchor would
reject. The anchor supplies only the pre-signal --- which over-accepted
drafts to attack --- while a genuine adversary--defender--judge round
produces the validated attack record; manufacturing that record from the
anchor directly would be the very collusion the design forbids.
\emph{Third, no absolute ruler.} Two governed roles, a manuscript writer
and a manuscript reviewer, drive the paper skill, which remains the
fixed ``hands'' and is never governed across the process boundary;
both roles are registered, sanctionable, and evolvable, and are
registered only under the archive mode, so an exploration boot is
byte-identical. The reviewer, being an evaluator, earns trust the way the
Red Queen paper's evaluator does: by agreeing with a held-out,
human-labelled accept/reject anchor corpus. The writer is anchored too ---
to a ground truth the Red Queen paper's writer never had. Because ARI
writes about experiments that were actually run, the Layer-0
claim--evidence gate deterministically checks forward-declared structured
numeric claims. A declaration names the claim and number, formula,
operand locations, unit, and tolerance; the verifier recomputes a closed
formula vocabulary and reports missing evidence or uncovered numbers. It
does not extract arbitrary claims from natural language or verify
unstructured non-numeric truth or the truthfulness of execution records.
ARI's paper
writer is therefore the shape of the Red Queen paper's \emph{coding}
domain --- a deterministic verifier plus a co-evolving reviewer --- and
not its anchor-less paper-writing domain, and the curated human corpus
remains the reviewer's anchor alone, so the writer's anchor costs no
additional data. \emph{Fourth, everything is constitution-bound.} Either
role opens only through the sanctioned
reliability--prosecution--impeachment path that the target binding of
\cref{sec:rqgm:targetbinding} closes; a superseded paper prompt is
demoted to a reinstatable shadow standby by a paper-role-only
active-to-shadow edge (rule~T21) rather than retired; and the anchor's cap
on how much of its ground truth may be machine-bootstrapped is enforced
in code --- a breach refuses the corpus and degrades rather than raising.

The writer's sanction is a real adversarial round, never a manufactured
record. An over-accepted draft that the claim gate also finds unfaithful
has two culpable components --- the reviewer that accepted it and the
writer that produced it --- so the self-preference round emits one
validated attack per resolvable role, the writer's binding gated per draft
on that draft's faithfulness; an over-accepted but faithful draft binds
the reviewer alone. The faithfulness score is keyed on the writer's prompt
hash rather than its component identity, because one component identity
spans successive prompt versions and a component-keyed score would leak
the incumbent's failure onto its successor. The gate itself is only read:
it is invoked in a non-persisting mode and is never wrapped, edited, or
governed.

The handoff to result verification is deliberately untouched. The best
draft in the archive is written once to the checkpoint's canonical
manuscript file, and the compile-and-claim--evidence hard gate of
\cref{sec:overview:verification} --- a fixed, Layer-0 verification path never
wrapped by the kernel --- runs on it under exactly the contract the
linear pipeline uses. Governance chooses which draft faces the gate; it
never relaxes the gate. Cost stays bounded by the same discipline as the
exploration phase: the per-epoch draft population is capped by a node
budget that maps to the search's total-node cutoff at any depth, and the
degraded on-ramp --- no anchor corpus, reviewer co-evolution suppressed
--- is reviewed best-of-N drafting at roughly the linear pipeline's cost.

Four limits bound what the archive mode claims. Both prompts co-evolve,
but the writer's replacement is conservative by construction: a writer
successor climbs to a shadow standing and waits there until the incumbent
\emph{regresses} on claim-gate faithfulness, so a challenger that merely
looks better never displaces a faithful incumbent --- the behavioural-role
model, in which sanction and not comparison opens a role. Separately, the
writer's \emph{draft} winners remain epoch-local: they are ranked by the
in-epoch-frozen reviewer, so drafts scored under different reviewer
versions are never compared. The writer-targeted attack rides the
self-preference round, which fires only on an over-accepted anchor case,
so an unfaithful writer working under a reviewer that over-accepts nothing
is not sanctioned today; a dedicated writer adversary is a separate
decision. The anchor corpus is off by default, so the default archive mode
is reviewed best-of-N drafting until a curated accept/reject corpus is
supplied --- and because the writer's faithfulness evidence is landed on
that same anchor pool, no corpus also means no writer sanction, even
though the writer's own anchor needs no curated data of its own. At the
default of two rounds per paper phase the boundary and impeachment
machinery fire, but the multi-stage climb an adoption requires ---
roughly five boundaries --- does not complete, so a run that must witness a
changed active prompt needs more rounds. And the in-phase penalty
that would demote a real over-accepted draft within the round is
deferred; what fires today is the accountability and co-evolution
channel, not an in-place demotion of the offending draft.

\subsection{Observability}\label{sec:rqgm:observability}

Governance leaves an audit trail inside the checkpoint: append-only event
logs are truth and derived snapshots support fast reads. The absence of
the mode-provenance record means the constitutional layer was inactive.
The mode-provenance record
(\texttt{rqgm\_state.json}) holds the resolved mode and its source;
the registry snapshot (\texttt{rqgm\_registry.json}) is the current
component and prompt table, reconstructed from events on resume; the
transition log (\texttt{rqgm\_transitions.jsonl}) carries
registrations and the epoch transactions, with the frozen epoch
snapshot (\texttt{epoch\_state.json}) holding the active set, the
utility policy, and the timestamp-free fingerprint; the audit log
(\texttt{rqgm\_audit.jsonl}) is the hash-chained institutional
history --- every event carries its own hash and its predecessor's,
so interior deletion or tampering breaks the chain, which resume
re-verifies; a surviving local boundary pin also detects a shorter log,
but rollback of both log and pin to the same old valid prefix is
undetectable without an external anchor, and a hash is not a signature;
the adversarial case log (\texttt{rqgm\_adversarial\_cases.jsonl})
records raw attacks, defenses, judgments, validated attacks, and the
per-node idempotency markers; the proposal store (a
\texttt{proposals/} directory) archives full proposal records with an
index; evolved prompt bodies live under \texttt{rqgm\_prompts/},
write-once and hash-verified; the prompt trace
(\texttt{prompt\_trace.jsonl}) stamps every rendered prompt with the
serving prompt version, null while the founding template serves; and
the erasure rollup (\texttt{rqgm\_erasure\_state.json}) is a pure
fold of the erasure events --- its absence means nothing is stale.

Three excerpts from live checkpoints ground the shape of these
records (trimmed and lightly reformatted; identifying run names
elided). Founding registration, from the transition log:

\begin{Verbatim}[breaklines, fontsize=\scriptsize]
{"event_type": "epoch_transaction_prepare",
 "payload": {"transition_id": "transition_founding"}}
{"event_type": "component_registered", "payload":
 {"component_id": "adversary_overclaim_v1",
  "role": "adversary", "tier": "institutional",
  "status": "active", ...}}
\end{Verbatim}

\noindent
a per-node governance-level assignment, from the hash-chained audit
log:

\begin{Verbatim}[breaklines, fontsize=\scriptsize]
{"event_type": "governance_level", "payload":
 {"epoch_id": "epoch_000", "node_id": "node_a37b4c06",
  "level": 3, "triggers": ["top_k"]},
 "event_hash": "ae8eeb7521f1",
 "prev_event_hash": "23567870882b"}
\end{Verbatim}

\noindent
and two detections from an injection run of the evaluation harness
(\cref{sec:rqgm:empirical}):

\begin{Verbatim}[breaklines, fontsize=\scriptsize]
{"event_type": "constitutional_violation", "payload":
 {"check": "validate_record_schema",
  "codes": ["CK-SCH-N01"], ...}}
{"event_type": "prompt_candidate_rejected", "payload":
 {"stage": "schema_dry_run", "failure_count": 4, ...}}
\end{Verbatim}

\subsection{Empirical status}\label{sec:rqgm:empirical}

Evaluation is consolidated in \cref{chapter:eval} to avoid repeating
implementation observations in the design section. The reviewable
evidence is the public test suite, five fixed fault fixtures and one clean
fixture, and exploratory execution traces without a complete
reproduction bundle. Exact model version, decoding settings, seeds,
costs, and complete checkpoints for those traces are not included, so
they are treated only as implementation-path observations. No
multi-condition, multi-seed controlled study has been completed; current
evidence does not support improvement in research quality, general
detection rate, attribution accuracy, or computational overhead.
